\scanpage{55}
\mspage{-1.5-}

product
\[
\mathbb B_1\times\mathbb B_2\times\cdots\times\mathbb B_m
\]
of affine spaces such that $G$ leaves each factor invariant and such that $\mathbb B_i$ is irreducible under the \unsure{actionof|action of} $G$. Suppose the factors are so numbered that $G$ acts trivially on $\mathbb B_1,\ldots,\mathbb B_t$ and non-trivially on $\mathbb B_{t+1},\ldots,\mathbb B_m$. Then $G$ acts trivially on
\[
\mathbb B_1\times\cdots\times\mathbb B_t\times(0)\times\cdots\times(0)\subset\mathbb A^n
\]
and this subspace contains every other subspace on which $G$ acts trivially. The projection of $\mathbb A^n$ onto this subspace which annihilates $\mathbb B_{t+1}\times\cdots\times\mathbb B_m$ is all-important in the sequel. Actually, it is more convenient on the dual level:

\uline{Definition 1.5}: Let $G$ be a reductive algebraic group, and let $S=\Gamma(G,\underline{o}_G)$. Let $\hat\sigma$ be a dual action of $G$ on a vector space $V$. Then a \uline{Reynolds operator} is a homomorphism $E:V\longrightarrow V$ such that:
\[
\begin{aligned}
\text{(i)}\quad &E\text{ commutes with }\hat\sigma,\text{ i.e. }\hat\sigma\circ E=(1_S\otimes E)\circ\hat\sigma.\\
\text{(ii)}\quad &E^2=E\\
\text{(iii)}\quad &Ex=x\text{ if and only if }\hat\sigma(x)=1\otimes x.
\end{aligned}
\]
If $V$ is finite-dimensional, it follows immediately from the definition of ``reductive'' that $E$ exists and is unique. The same holds for a general $V$ by virtue of the above lemma, (compare Cartier [6], exposé 7). Because of the canonical nature of $E$, whenever we have dual actions of $G$ on two vector spaces $V_1$ and $V_2$, and a linear map

\margin{54}
